% Part: first-order-logic
% Chapter: models-theories
% Section: size-of-structures

\documentclass[../../../include/open-logic-section]{subfiles}

\begin{document}

\olfileid{fol}{mat}{siz}

\olsection{\printtoken{P}{structure} चे आकारमान व्यक्त करणे}

\begin{explain}
भाषेतील अतार्किक चिन्हे न वापरताही रचनांचे काही गुणधर्म व्यक्त करता येतात.
उदाहरणार्थ, अशी !!{sentence}s आहेत की ती !!a{structure} मध्ये तेव्हा आणि केवळ
तेव्हाच सत्य असतात, जेव्हा त्या !!{structure} च्या !!{domain} मध्ये
किमान, जास्तीत जास्त किंवा नेमके दिलेल्या संख्येइतके~$n$ !!{element}s असतात.
\end{explain}

\begin{prop}
!!{sentence}
\begin{multline*}
  !A_{\ge n} \ident \lexists[x_1][\lexists[x_2][\dots\lexists[x_n][{}]]]\\
  \begin{aligned}
  (\eq/[x_1][x_2] \land {}
  \eq/[x_1][x_3] \land \eq/[x_1][x_4] \land \dots \land \eq/[x_1][x_n] \land {}\\
\eq/[x_2][x_3] \land \eq/[x_2][x_4] \land \dots \land {} \eq/[x_2][x_n] \land {} \\
\vdots\\
\eq/[x_{n-1}][x_n])
\end{aligned}
\end{multline*}
हे !!a{structure}~$\Struct M$ मध्ये तेव्हा आणि केवळ तेव्हाच सत्य असते,
जेव्हा $\Domain M$ मध्ये किमान $n$ !!{element}s असतात. परिणामी,
$\Sat{M}{\lnot !A_{\ge n+1}}$ तेव्हा आणि केवळ तेव्हाच लागू होते, जेव्हा
$\Domain M$ मध्ये जास्तीत जास्त~$n$ !!{element}s असतात.

\end{prop}

\begin{prop}
!!{sentence}
\begin{multline*}
!A_{= n} \ident \lexists[x_1][\lexists[x_2][\dots\lexists[x_n][{}]]] \\
\begin{aligned}
  (\eq/[x_1][x_2] \land {}
  \eq/[x_1][x_3] \land \eq/[x_1][x_4] \land \dots \land \eq/[x_1][x_n] \land {}\\
\eq/[x_2][x_3] \land \eq/[x_2][x_4] \land \dots \land {} \eq/[x_2][x_n] \land {} \\
\vdots\\
\eq/[x_{n-1}][x_n] \land {} \\
\lforall[y][(\eq[y][x_1] \lor \dots \lor \eq[y][x_n]]))
\end{aligned}
\end{multline*}
हे !!a{structure}~$\Struct M$ मध्ये तेव्हा आणि केवळ तेव्हाच सत्य असते,
जेव्हा $\Domain M$ मध्ये नेमके $n$ !!{element}s असतात.
\end{prop}

\begin{prop}
एखादी !!{structure} अनंत तेव्हा आणि केवळ तेव्हाच असते, जेव्हा ती पुढील
संचाचे प्रतिमान असते:
\[
\{!A_{\ge 1}, !A_{\ge 2}, !A_{\ge 3}, \dots \}.
\]
\end{prop}

असे एकही पूर्णतः तार्किक वाक्य नाही की जे~$\Struct M$ मध्ये तेव्हा आणि केवळ
तेव्हाच सत्य असेल, जेव्हा $\Domain M$ अनंत असेल. पण अशी अतार्किक
!!{predicate}s असलेली !!{sentence}s देता येतात ज्यांची फक्त अनंत प्रतिमाने
असतात (जरी प्रत्येक अनंत !!{structure} त्यांचे प्रतिमान नसते). सांत रचना
असण्याचा गुणधर्म आणि !!{nonenumerable} रचना असण्याचा गुणधर्म
!!{sentence}s च्या अनंत संचानेही व्यक्त करता येत नाही. ही तथ्ये संहतता आणि
लोव्हेनहाइम--स्कोलेम प्रमेयांवरून निष्पन्न होतात.

\end{document}
